\documentclass[11pt]{article}
\usepackage{mathblatt}
\def\mitzone{1}
\begin{document}
\blattkopf{Matrizen und Übergangsprozesse}{Gesamt}
\verzeichniszeile{\verz{zone}{Kennst du schon (Nr.~1–8)} \verztrenn \verz{e1}{1 Matrizen als Rechenobjekte (Nr.~9–10)} \verztrenn \verz{e2}{2 Vektoren unter Matrizen (Nr.~11–12)} \verztrenn \verz{e3}{3 Verflechtung (Nr.~13–14)} \verztrenn \verz{e4}{4 Übergangsmodell (Nr.~15–16)} \verztrenn \verz{e5}{5 Stationär und langfristig (Nr.~17–19)} \verztrenn \verz{abhaken}{Das kann ich}}
% Zone „Das kennst du schon“ – aus bank/matrizen-und-uebergangsprozesse/zone.jsonl (zusammenbau v0.1)
\einheitenkopf[zone][Kennst du schon]{Das kennst du schon}
\setcounter{aufgabe}{0}

%% TODO Titel: Ich-kann-Satz fehlt in der Bank (Platzhalter: Kettenname) – Nr. 1
%% TODO Anweisung: ein Satz über der ersten Teilaufgabe fehlt in der Bank – Nr. 1
\begin{aufgabe}{Lineare Gleichungssysteme mit zwei und drei Unbekannten lösen, auch unterbestimmte (Parameter frei wählen) und überbestimmte (überzählige Gleichung als Probe)}
\begin{gleichungsraster}[2]
\gl{\text{Löse das Gleichungssystem: $x + y = 7$ und $x - y = 3$.}} &
\gl{\text{Löse das Gleichungssystem: $0{,}4x + 0{,}5y = 7$ und $x + y = 16$.}} \\
\end{gleichungsraster}
\end{aufgabe}

%% TODO Titel: Ich-kann-Satz fehlt in der Bank (Platzhalter: Kettenname) – Nr. 2
%% TODO Anweisung: ein Satz über der ersten Teilaufgabe fehlt in der Bank – Nr. 2
\begin{aufgabe}{Vektoren als Listen lesen und schreiben (Zustands-, Mengen-, Kostenvektoren), Komponenten adressieren}
\begin{teile}
\teil Ein Kiosk verkauft am Montag 40 Brezeln, 25 Brötchen und 12 Croissants. Schreibe die Mengen als Vektor (Brezeln $\mid$ Brötchen $\mid$ Croissants). ( \leerfeld | \leerfeld | \leerfeld )
\teil Berechne $2 \cdot (1{,}5 \mid 0{,}4 \mid 3) + (0{,}5 \mid 1{,}2 \mid 1)$. ( \leerfeld | \leerfeld | \leerfeld )
\end{teile}
\end{aufgabe}

%% TODO Titel: Ich-kann-Satz fehlt in der Bank (Platzhalter: Kettenname) – Nr. 3
%% TODO Anweisung: ein Satz über der ersten Teilaufgabe fehlt in der Bank – Nr. 3
\begin{aufgabe}{Diagramme lesen und zeichnen (Knoten, Pfeile, Punktdiagramme, Kurvenformen unterscheiden)}
\begin{teile}
\teil Lies ab: Wie viele Gäste kamen am Dienstag?

\saeulenab[ymax=100,ystep=20,yfein=5,ylabel=Gäste]{Mo/35,Di/60,Mi/45}
\leerfeld Gäste
\teil Lies ab und entscheide: Um welchen Faktor wächst die Anzahl von Woche zu Woche? Ist das Wachstum linear oder exponentiell?

\liniendia[ymax=40,ystep=10,yfein=5,ylabel=Anzahl]{W1/5,W2/10,W3/20,W4/40}
Faktor \leerfeld
\end{teile}
\end{aufgabe}

%% TODO Titel: Ich-kann-Satz fehlt in der Bank (Platzhalter: Kettenname) – Nr. 4
%% TODO Anweisung: ein Satz über der ersten Teilaufgabe fehlt in der Bank – Nr. 4
\begin{aufgabe}{Terme mit einer Unbekannten aufstellen und gegen Schranken auswerten (linear, mit Randwerten)}
\begin{teile}
\teil Eine Taxifahrt kostet 4\,€ Grundpreis und 2\,€ je Kilometer. Wie viele Kilometer kannst du mit 20\,€ höchstens fahren? \leerfeld[km]
\teil Der Anteil $0{,}2 + 0{,}5p$ hängt von p ab, p liegt zwischen 0 und 1. Zwischen welchen Werten liegt der Anteil? Setze die Randwerte ein. zwischen \leerfeld und \leerfeld
\end{teile}
\end{aufgabe}

%% TODO Titel: Ich-kann-Satz fehlt in der Bank (Platzhalter: Kettenname) – Nr. 5
%% TODO Anweisung: ein Satz über der ersten Teilaufgabe fehlt in der Bank – Nr. 5
\begin{aufgabe}{Anteile, Prozentsätze und Anzahlen ineinander umrechnen (Bleibe- und Wechselanteile, Gegenanteil)}
\begin{teile}
\teil Von 200 Kunden wechseln 10\,\% zu einem anderen Geschäft. Wie viele wechseln? \leerfeld Kunden
\teil Von 450 Fahrgästen bleiben 84\,\% im Bus sitzen. Wie viele bleiben sitzen? \leerfeld Fahrgäste
\end{teile}
\end{aufgabe}

%% TODO Titel: Ich-kann-Satz fehlt in der Bank (Platzhalter: Kettenname) – Nr. 6
%% TODO Anweisung: ein Satz über der ersten Teilaufgabe fehlt in der Bank – Nr. 6
\begin{aufgabe}{Potenzen mit Basis unter eins fallen, Exponentialungleichungen durch Probieren oder Logarithmieren lösen, Ergebnisse sinnvoll aufrunden}
\begin{teile}
\teil Berechne $0{,}5^3$. \leerfeld
\teil Für welche kleinste natürliche Zahl n gilt $1{,}2^n > 2$? Probiere. n = \leerfeld
\end{teile}
\end{aufgabe}

%% TODO Titel: Ich-kann-Satz fehlt in der Bank (Platzhalter: Kettenname) – Nr. 7
%% TODO Anweisung: ein Satz über der ersten Teilaufgabe fehlt in der Bank – Nr. 7
\begin{aufgabe}{Lineare Gleichungssysteme mit zwei und drei Unbekannten lösen, auch unterbestimmte (Parameter frei wählen) und überbestimmte (überzählige Gleichung als Probe) – Fehler finden}
\begin{teile}
\teil Tom soll alle Lösungen von $3x + y = 12$ angeben. Er schreibt: „$x = 4$, $y = 0$ – das ist die Lösung.“ Finde den Fehler und gib alle Lösungen an.
\end{teile}
\end{aufgabe}

%% TODO Titel: Ich-kann-Satz fehlt in der Bank (Platzhalter: Kettenname) – Nr. 8
%% TODO Anweisung: ein Satz über der ersten Teilaufgabe fehlt in der Bank – Nr. 8
\begin{aufgabe}{Lineare Gleichungssysteme mit zwei und drei Unbekannten lösen, auch unterbestimmte (Parameter frei wählen) und überbestimmte (überzählige Gleichung als Probe) – selbst rechnen}
\begin{gleichungsraster}[2]
\gl{\text{Gib alle Lösungen von $2x + y = 10$ an (setze $x = t$). Nenne die Lösung für $t = 3$.}} & \\
\end{gleichungsraster}
\end{aufgabe}

\clearpage
% Einheit 1 von 5 – Matrizen als Rechenobjekte (bank/matrizen-und-uebergangsprozesse/e1.jsonl, zusammenbau v0.1)
%% TODO Zeitmarke Einheit 1: Marken-Zeile nicht lesbar
%% TODO Zweigzeile Teil 1: was hier gelernt wird (Satz in Schülersprache aus dem Einheitstitel) – Einheit 1
\einheitenkopf[e1]{Einheit 1 von 5 · Matrizen als Rechenobjekte}
\zweigzeile{keine Prüfungsaufgabe}
\setcounter{aufgabe}{8}

%% TODO Titel: Ich-kann-Satz fehlt in der Bank (Platzhalter: Kettenname) – Nr. 9
%% TODO Anweisung: ein Satz über der ersten Teilaufgabe fehlt in der Bank – Nr. 9
\begin{aufgabe}{Matrizen als Rechenobjekte}
%% TODO \rechenplatz{n}: Schrittzahl des Grundfalls fehlt in der Bank (nur bei fester Schreibform, 2.3 b)
\begin{teile}
\teil Rechnen oder deuten? Nur ankreuzen, nicht ausführen: „Ermittle das Produkt der Matrix $((2 | 5), (1 | 4))$ mit dem Vektor $(3 | 2)$.“ \\ \kreuz{rechnen} \\ \kreuz{deuten}
\teil Berechne $A \cdot v$ für $A = ((2 | 1), (3 | 4))$ und $v = (1 | 2)$. ( \leerfeld | \leerfeld )
\teil A hat 2 Zeilen und 3 Spalten, B hat 3 Zeilen und 4 Spalten. Ist $A \cdot B$ bildbar, und welches Format hat es? Ist $B \cdot A$ bildbar?
\teil Berechne $M^2$ für $M = ((1 | 2), (2 | 1))$. (( \leerfeld | \leerfeld ), ( \leerfeld | \leerfeld ))
\teil Bestimme die Inverse von $A = ((2 | 1), (1 | 1))$ über den Ansatz $A \cdot B = E$ mit $B = ((a | b), (c | d))$. (( \leerfeld | \leerfeld ), ( \leerfeld | \leerfeld ))
\teil Beschreibe die Wirkung von $P = ((0 | 1), (1 | 0))$ auf einen Vektor $(a | b)$.
\end{teile}
\begin{gleichungsraster}[2]
\gl{\text{$M = ((a | 1), (2 | b))$ und $M \cdot (1 | 2) = (5 | 10)$. Bestimme a und b.}} &
\end{gleichungsraster}
\begin{teile}
\teil $A = ((0 | 0), (1 | 0))$ und $X = ((a | b), (c | d))$ mit reellen Zahlen a, b, c, d. Ermittle alle a, b, c, d, für die $(X + A)^2 = X^2 + 2 \cdot A \cdot X + A^2$ gilt \hfill (Abitur 2025 LK).
\end{teile}
\end{aufgabe}

%% TODO Titel: Ich-kann-Satz fehlt in der Bank (Platzhalter: Kettenname) – Nr. 10
%% TODO Anweisung: ein Satz über der ersten Teilaufgabe fehlt in der Bank – Nr. 10
\begin{aufgabe}{Matrizen als Rechenobjekte – Fehler finden, Begründen}
\begin{teile}
\teil Lena berechnet das Quadrat von $M = ((2 | 1), (0 | 3))$ und schreibt $M^2 = ((4 | 1), (0 | 9))$. Finde den Fehler und rechne richtig.
\teil Begründe ohne Rechnung, warum $A \cdot B$ und $B \cdot A$ bei Matrizen verschieden sein können.
\end{teile}
\end{aufgabe}

\clearpage
% Einheit 2 von 5 – Vektoren unter Matrizen (bank/matrizen-und-uebergangsprozesse/e2.jsonl, zusammenbau v0.1)
%% TODO Zeitmarke Einheit 2: Marken-Zeile nicht lesbar
%% TODO Zweigzeile Teil 1: was hier gelernt wird (Satz in Schülersprache aus dem Einheitstitel) – Einheit 2
\einheitenkopf[e2]{Einheit 2 von 5 · Vektoren unter Matrizen}
\zweigzeile{keine Prüfungsaufgabe}
\setcounter{aufgabe}{10}

%% TODO Titel: Ich-kann-Satz fehlt in der Bank (Platzhalter: Kettenname) – Nr. 11
%% TODO Anweisung: ein Satz über der ersten Teilaufgabe fehlt in der Bank – Nr. 11
\begin{aufgabe}{Vektoren unter Matrizen}
%% TODO \rechenplatz{n}: Schrittzahl des Grundfalls fehlt in der Bank (nur bei fester Schreibform, 2.3 b)
\begin{teile}
\teil Allgemein oder am Beispiel? Nur ankreuzen: „Zeige, dass jede Matrix $((p | q), (r | s))$ mit $p + r = 1$ und $q + s = 1$ die Komponentensumme jedes Vektors erhält.“ \\ \kreuz{allgemein zeigen} \\ \kreuz{ein Beispiel genügt}
\end{teile}
\begin{gleichungsraster}[2]
\gl{\text{Bestimme alle Vektoren v mit $M \cdot v = v$ für $M = ((0{,}6 | 0{,}2), (0{,}4 | 0{,}8))$.}} &
\gl{\text{Bestimme alle Vektoren v mit $M \cdot v = 3 \cdot v$ für $M = ((2 | 1), (0 | 3))$.}} \\
\gl{\text{$M = ((1 | 2 | 0), (0 | 1 | 1), (2 | 0 | 1))$ und $M \cdot (a | 1 | 2) = (5 | 3 | 8)$. Bestimme a und prüfe mit den übrigen Zeilen.}} &
\end{gleichungsraster}
\begin{teile}
\teil Zeige allgemein: Ist $M = ((p | q), (r | s))$ mit $p + r = 1$ und $q + s = 1$, so hat $M \cdot (x | y)$ dieselbe Komponentensumme wie $(x | y)$.
\end{teile}
\begin{gleichungsraster}[2]
\gl{\text{Eine Matrix heißt selbstinvers, wenn $M \cdot M = E$ gilt. Für welche a ist $M = ((a | 1), (0 | -1))$ selbstinvers?}} &
\end{gleichungsraster}
\begin{teile}
\teil $M = ((1 | 2), (2 | 4))$. Zeige, dass $M \cdot v = (0 | 0)$ außer dem Nullvektor weitere Lösungen hat, und gib alle an.
\teil $A = ((3 | 2), (1 | 4))$ und $w = (1 | 1)$ mit $A \cdot w = 5 \cdot w$. Für jede Zahl c gibt es eine Zahl b mit $A^{-1} \cdot (b \cdot w) = c \cdot w$. Bestimme b in Abhängigkeit von c, ohne $A^{-1}$ zu berechnen \hfill (Abitur 2026 LK).
\end{teile}
\end{aufgabe}

%% TODO Titel: Ich-kann-Satz fehlt in der Bank (Platzhalter: Kettenname) – Nr. 12
%% TODO Anweisung: ein Satz über der ersten Teilaufgabe fehlt in der Bank – Nr. 12
\begin{aufgabe}{Vektoren unter Matrizen – Fehler finden, Begründen}
\begin{teile}
\teil Jan soll alle v mit $M \cdot v = v$ für $M = ((0{,}5 | 0{,}25), (0{,}5 | 0{,}75))$ bestimmen. Er schreibt: „$v = (0 | 0)$, denn der Nullvektor erfüllt die Gleichung.“ Finde den Fehler und gib alle Lösungen an.
\teil Begründe: Erfüllt v die Gleichung $M \cdot v = t \cdot v$, dann erfüllt auch $5v$ sie.
\end{teile}
\end{aufgabe}

\clearpage
% Einheit 3 von 5 – Verflechtung (bank/matrizen-und-uebergangsprozesse/e3.jsonl, zusammenbau v0.1)
%% TODO Zeitmarke Einheit 3: Marken-Zeile nicht lesbar
%% TODO Zweigzeile Teil 1: was hier gelernt wird (Satz in Schülersprache aus dem Einheitstitel) – Einheit 3
\einheitenkopf[e3]{Einheit 3 von 5 · Verflechtung}
\zweigzeile{keine Prüfungsaufgabe}
\setcounter{aufgabe}{12}

%% TODO Titel: Ich-kann-Satz fehlt in der Bank (Platzhalter: Kettenname) – Nr. 13
%% TODO Anweisung: ein Satz über der ersten Teilaufgabe fehlt in der Bank – Nr. 13
\begin{aufgabe}{Verflechtung}
%% TODO \rechenplatz{n}: Schrittzahl des Grundfalls fehlt in der Bank (nur bei fester Schreibform, 2.3 b)
\begin{teile}
\teil Die Bedarfsmatrix A hat die Zeilen R1, R2 (Rohstoffe) und die Spalten Z1, Z2 (Zwischenprodukte). In Zeile 1, Spalte 2 steht 4. Welcher Pfeil im Verflechtungsdiagramm trägt die 4? \\ \kreuz{von R1 nach Z2} \\ \kreuz{von R2 nach Z1}
\teil Eine Bäckerei braucht je Blech Apfelkuchen 2 kg Mehl und 1 kg Zucker, je Blech Nusskuchen 3 kg Mehl und 2 kg Zucker: $A = ((2 | 3), (1 | 2))$ (Zeilen Mehl, Zucker). Berechne den Bedarf für 5 Bleche Apfel- und 4 Bleche Nusskuchen. ( \leerfeld | \leerfeld )
\teil Im Verflechtungsdiagramm führen die Pfeile R1 → Z1 mit 2, R1 → Z2 mit 5, R2 → Z2 mit 1 und R3 → Z1 mit 4. Gib die Matrix A an (Zeilen R1, R2, R3; Spalten Z1, Z2). (( \leerfeld | \leerfeld ), ( \leerfeld | \leerfeld ), ( \leerfeld | \leerfeld ))
\teil Ein Betrieb verarbeitet 4 Rohstoffe zu 3 Zwischenprodukten und daraus zu 2 Endprodukten. Welches Format haben A (Rohstoffe je Zwischenprodukt), B (Zwischenprodukte je Endprodukt) und $C = A \cdot B$?
\teil $A = ((1 | 2), (3 | 0))$ (Zeilen R1, R2; Spalten Z1, Z2) und $B = ((2 | 1), (1 | 4))$ (Zeilen Z1, Z2; Spalten E1, E2). Bestätige über die Pfade den Eintrag von $C = A \cdot B$ in Zeile 1, Spalte 2.
\end{teile}
\begin{gleichungsraster}[2]
\gl{\text{$C = ((2 | 1), (1 | 4))$ (Zeilen R1, R2; Spalten E1, E2). Verbraucht wurden 50 ME R1 und 95 ME R2. Wie viele Stück E1 und E2 wurden hergestellt?}} &
\end{gleichungsraster}
\begin{teile}
\teil Je Stück eines Produkts braucht man 3 kg R1, 5 kg R2 und 2 kg R3. Im Lager sind 120 kg R1, 150 kg R2 und 90 kg R3. Wie viele Stück kann man höchstens herstellen? \leerfeld[Stück]
\teil Eine Möbelwerkstatt fertigt aus Gestellen (Z1) und Platten (Z2) Tische (E1) und Bänke (E2); $B = ((2 | 1), (0 | 3))$ gibt die Zwischenprodukte je Endprodukt an. Für einen Auftrag über 20 Tische und 10 Bänke gilt $(z | 3z) \cdot B \cdot (20 | 10) = 1400$. Erläutere die Gleichung im Sachzusammenhang und löse sie \hfill (Abitur 2025 LK).
\end{teile}
\end{aufgabe}

%% TODO Titel: Ich-kann-Satz fehlt in der Bank (Platzhalter: Kettenname) – Nr. 14
%% TODO Anweisung: ein Satz über der ersten Teilaufgabe fehlt in der Bank – Nr. 14
\begin{aufgabe}{Verflechtung – Fehler finden, Begründen, Anwendung}
\begin{teile}
\teil $A = ((1 | 2), (2 | 1))$ gibt Rohstoffe je Zwischenprodukt, $B = ((1 | 1), (0 | 2))$ Zwischenprodukte je Endprodukt an. Für 10 Stück E1 und 20 Stück E2 rechnet Kai den Rohstoffbedarf als $A \cdot (10 | 20) = (50 | 40)$. Finde den Fehler und rechne richtig.
\teil Begründe, warum die Gesamtmatrix das Produkt $A \cdot B$ ist und nicht $B \cdot A$ (A: Rohstoffe je Zwischenprodukt, B: Zwischenprodukte je Endprodukt).
\teil Eine Kelterei mischt aus Apfel- und Traubensaft zwei Säfte: je Liter „Mild“ $0{,}7$ l Apfel und $0{,}3$ l Traube, je Liter „Herb“ $0{,}4$ l Apfel und $0{,}6$ l Traube. Im Lager sind 1000 l Apfel- und 800 l Traubensaft. Kann die Kelterei einen Auftrag über 1000 l Mild und 500 l Herb erfüllen?
\end{teile}
\end{aufgabe}

\clearpage
% Einheit 4 von 5 – Übergangsmodell (bank/matrizen-und-uebergangsprozesse/e4.jsonl, zusammenbau v0.1)
%% TODO Zeitmarke Einheit 4: Marken-Zeile nicht lesbar
%% TODO Zweigzeile Teil 1: was hier gelernt wird (Satz in Schülersprache aus dem Einheitstitel) – Einheit 4
\einheitenkopf[e4]{Einheit 4 von 5 · Übergangsmodell}
\zweigzeile{keine Prüfungsaufgabe}
\setcounter{aufgabe}{14}

%% TODO Titel: Ich-kann-Satz fehlt in der Bank (Platzhalter: Kettenname) – Nr. 15
%% TODO Anweisung: ein Satz über der ersten Teilaufgabe fehlt in der Bank – Nr. 15
\begin{aufgabe}{Übergangsmodell}
%% TODO \rechenplatz{n}: Schrittzahl des Grundfalls fehlt in der Bank (nur bei fester Schreibform, 2.3 b)
\begin{teile}
\teil Es gilt $v_{n+1} = M \cdot v_n$; die Zustände sind A, B, C in dieser Reihenfolge. In Zeile 3, Spalte 1 von M steht $0{,}2$. Von wo nach wo wechseln diese 20\,\%? \\ \kreuz{von A nach C} \\ \kreuz{von C nach A}
\teil Es gilt $v_{n+1} = M \cdot v_n$; Zustände: Studio A, Studio B; $M = ((0{,}9 | 0{,}15), (0{,}1 | 0{,}85))$. Zeichne das Übergangsdiagramm.

\rechenplatz{4}
\teil Übergangsdiagramm: Schleife an A mit $0{,}8$, Pfeil A → B mit $0{,}2$, Pfeil B → A mit $0{,}35$, Schleife an B mit $0{,}65$. Stelle M auf (Reihenfolge A, B; $v_{n+1} = M \cdot v_n$). (( \leerfeld | \leerfeld ), ( \leerfeld | \leerfeld ))
\teil Es gilt $v_{n+1} = M \cdot v_n$; Studierende wählen täglich zwischen den Mensen P und Q (Reihenfolge P, Q); $M = ((0{,}75 | 0{,}4), (0{,}25 | 0{,}6))$. Deute den Eintrag $0{,}4$.
\teil Es gilt $v_{n+1} = M \cdot v_n$; $M = ((0{,}8 | 0{,}3), (0{,}2 | 0{,}7))$, $v_0 = (500 | 500)$. Berechne $v_1$ und $v_2$. v1 = ( \leerfeld | \leerfeld ), v2 = ( \leerfeld | \leerfeld )
\end{teile}
\begin{gleichungsraster}[2]
\gl{\text{Es gilt $v_{n+1} = M \cdot v_n$; $M = ((0{,}7 | 0{,}1), (0{,}3 | 0{,}9))$. Jetzt ist die Verteilung $(340 | 660)$. Bestimme die Verteilung einen Schritt vorher.}} &
\gl{\text{Es gilt $v_{n+1} = M \cdot v_n$; $M = ((0{,}8 | 0{,}1), (0{,}2 | 0{,}9))$, $v_0 = (a | 600)$. Nach einem Schritt sind 460 bei A. Bestimme a.}} \\
\end{gleichungsraster}
\begin{teile}
\teil Es gilt $v_{n+1} = M \cdot v_n$; zwei Buchläden A, B teilen sich 1000 Stammkunden; $M = ((0{,}9 | 0{,}2), (0{,}1 | 0{,}8))$. Begründe, dass $(165 | 835)$ im Modell nicht die Kundenverteilung eines Jahres sein kann, das auf ein anderes Modelljahr folgt \hfill (Abitur 2026 GK).
\end{teile}
\end{aufgabe}

%% TODO Titel: Ich-kann-Satz fehlt in der Bank (Platzhalter: Kettenname) – Nr. 16
%% TODO Anweisung: ein Satz über der ersten Teilaufgabe fehlt in der Bank – Nr. 16
\begin{aufgabe}{Übergangsmodell – Fehler finden, Begründen, Anwendung}
\begin{teile}
\teil Übergangsdiagramm: A bleibt $0{,}6$, A → B $0{,}4$, B → A $0{,}1$, B bleibt $0{,}9$. Emil schreibt $M = ((0{,}6 | 0{,}4), (0{,}1 | 0{,}9))$ für $v_{n+1} = M \cdot v_n$. Finde den Fehler und gib M richtig an.
\teil Begründe, warum in $v_{n+1} = M \cdot v_n$ die Spalten von M zu den Ausgangszuständen gehören.
\teil Ein Betreiber hat zwei Fitnessstudios. Von A bleiben monatlich 90\,\% der Mitglieder, 10\,\% wechseln nach B; von B bleiben 80\,\%, 20\,\% wechseln nach A. Jetzt hat A 1200 und B 800 Mitglieder. Studio B braucht mindestens 700 Mitglieder. Reicht das in den nächsten zwei Monaten?
\end{teile}
\end{aufgabe}

\clearpage
% Einheit 5 von 5 – Stationär und langfristig (bank/matrizen-und-uebergangsprozesse/e5.jsonl, zusammenbau v0.1)
%% TODO Zeitmarke Einheit 5: Marken-Zeile nicht lesbar
%% TODO Zweigzeile Teil 1: was hier gelernt wird (Satz in Schülersprache aus dem Einheitstitel) – Einheit 5
\einheitenkopf[e5]{Einheit 5 von 5 · Stationär und langfristig}
\zweigzeile{keine Prüfungsaufgabe}
\setcounter{aufgabe}{16}

%% TODO Titel: Ich-kann-Satz fehlt in der Bank (Platzhalter: Kettenname) – Nr. 17
%% TODO Anweisung: ein Satz über der ersten Teilaufgabe fehlt in der Bank – Nr. 17
\begin{aufgabe}{Stationär und langfristig}
%% TODO \rechenplatz{n}: Schrittzahl des Grundfalls fehlt in der Bank (nur bei fester Schreibform, 2.3 b)
\begin{teile}
\teil Es gilt $v_{n+1} = M \cdot v_n$; v ist die heutige Verteilung. Was beschreibt $M^3 \cdot v$? \\ \kreuz{drei Schritte vorwärts} \\ \kreuz{drei Schritte rückwärts} \\ \kreuz{einen Schritt vorwärts}
\end{teile}
\begin{gleichungsraster}[2]
\gl{\text{Es gilt $v_{n+1} = M \cdot v_n$; $M = ((0{,}8 | 0{,}4), (0{,}2 | 0{,}6))$, Gesamtzahl 900. Berechne die stationäre Verteilung.}} &
\end{gleichungsraster}
\begin{teile}
\teil Es gilt $v_{n+1} = M \cdot v_n$; $M = ((0{,}6 | 0), (0{,}4 | 1))$ (Reihenfolge A, B), Gesamtzahl 500. Gib die stationäre Verteilung an. ( \leerfeld | \leerfeld )
\end{teile}
\begin{gleichungsraster}[2]
\gl{\text{Es gilt $v_{n+1} = M \cdot v_n$; $M = ((0{,}8 | p), (0{,}2 | 1 - p))$; die Verteilung $(300 | 200)$ ist stationär. Bestimme p.}} &
\end{gleichungsraster}
\begin{teile}
\teil Es gilt $v_{n+1} = M \cdot v_n$; $M = ((0{,}9 | 0{,}2), (0{,}1 | 0{,}8))$ (Reihenfolge A, B), $v = (800 | 400)$. Wie viele wechseln von A nach B, wie viele von B nach A? Was folgt daraus?
\teil Es gilt $v_{n+1} = M \cdot v_n$; $M = ((0 | 0 | 4), (0{,}5 | 0 | 0), (0 | 0{,}5 | 0))$. Zeige $M^3 = E$ und deute das Ergebnis.
\teil Es gilt $v_{n+1} = M \cdot v_n$; $M = ((0 | 3), (0{,}5 | 0{,}5))$ und $v = (6 | 3)$. Zeige, dass $M \cdot v$ ein Vielfaches von v ist, und gib den Faktor an.
\teil Eine Lichterkette wechselt jede Sekunde zwischen Rot, Gelb und Blau (in dieser Reihenfolge) nach $K_b = ((b | 2b | 0), (2b | b | 0), (1 - 3b | 1 - 3b | 1))$ mit $0 \le b \le \frac{1}{3}$. Für $b < \frac{1}{3}$ nähern sich die Potenzen $K_b^k$ der Matrix mit den Zeilen $(0 | 0 | 0)$, $(0 | 0 | 0)$, $(1 | 1 | 1)$. Beschreibe die Farbwechsel in Abhängigkeit von b \hfill (Abitur 2022 LK).
\end{teile}
\end{aufgabe}

%% TODO Titel: Ich-kann-Satz fehlt in der Bank (Platzhalter: Kettenname) – Nr. 18
%% TODO Anweisung: ein Satz über der ersten Teilaufgabe fehlt in der Bank – Nr. 18
\begin{aufgabe}{Übergangsprozess: Eignung eines Populationsmodells zur langfristigen Beschreibung beurteilen}
\begin{teile}
\teil Eine Vogelpopulation auf einer kleinen Insel wächst im Modell jedes Jahr um den Faktor $1{,}2$. Beurteile, ob das Modell die Entwicklung über 50 Jahre beschreiben kann.
\end{teile}
\end{aufgabe}

%% TODO Titel: Ich-kann-Satz fehlt in der Bank (Platzhalter: Kettenname) – Nr. 19
%% TODO Anweisung: ein Satz über der ersten Teilaufgabe fehlt in der Bank – Nr. 19
\begin{aufgabe}{Stationär und langfristig – Fehler finden, Begründen, Anwendung}
\begin{teile}
\teil Es gilt $v_{n+1} = M \cdot v_n$; $M = ((0{,}9 | 0{,}4), (0{,}1 | 0{,}6))$, der Verein hat 1000 Mitglieder in zwei Gruppen. Anna sagt: „Stationär ist $(4 | 1)$.“ Finde den Fehler und gib die stationäre Verteilung an.
\teil Begründe, warum man beim Fixvektor einer Übergangsmatrix mit zwei Zuständen eine der beiden Gleichungen durch die Gesamtzahl ersetzt.
\teil In einem Ort gibt es zwei Bäckereien. Von A wechseln jede Woche 10\,\% der Kunden zu B, von B 15\,\% zu A; zusammen sind es 2000 Kunden. Bäckerei A braucht ab 1100 Kunden eine zweite Verkäuferin. Braucht sie auf Dauer eine?
\end{teile}
\end{aufgabe}

% Abhakseite „Das kann ich“ – Zone nur im Gesamt (\mitzone)
%% TODO Ich-kann-Sätze: Platzhalter wie in den Aufgabendateien
\begin{abhakseite}
\ifdefined\mitzone
\abhakgruppe{Kennst du schon}
\abhak{1}{Lineare Gleichungssysteme mit zwei und drei Unbekannten lösen, auch unterbestimmte (Parameter frei wählen) und überbestimmte (überzählige Gleichung als Probe)}
\abhak{2}{Vektoren als Listen lesen und schreiben (Zustands-, Mengen-, Kostenvektoren), Komponenten adressieren}
\abhak{3}{Diagramme lesen und zeichnen (Knoten, Pfeile, Punktdiagramme, Kurvenformen unterscheiden)}
\abhak{4}{Terme mit einer Unbekannten aufstellen und gegen Schranken auswerten (linear, mit Randwerten)}
\abhak{5}{Anteile, Prozentsätze und Anzahlen ineinander umrechnen (Bleibe- und Wechselanteile, Gegenanteil)}
\abhak{6}{Potenzen mit Basis unter eins fallen, Exponentialungleichungen durch Probieren oder Logarithmieren lösen, Ergebnisse sinnvoll aufrunden}
\abhak{7}{Lineare Gleichungssysteme mit zwei und drei Unbekannten lösen, auch unterbestimmte (Parameter frei wählen) und überbestimmte (überzählige Gleichung als Probe) – Fehler finden}
\abhak{8}{Lineare Gleichungssysteme mit zwei und drei Unbekannten lösen, auch unterbestimmte (Parameter frei wählen) und überbestimmte (überzählige Gleichung als Probe) – selbst rechnen}
\fi
\abhakgruppe{Einheit 1 · Matrizen als Rechenobjekte}
\abhak{9}{Matrizen als Rechenobjekte}
\abhak{10}{Matrizen als Rechenobjekte – Fehler finden, Begründen}
\abhakgruppe{Einheit 2 · Vektoren unter Matrizen}
\abhak{11}{Vektoren unter Matrizen}
\abhak{12}{Vektoren unter Matrizen – Fehler finden, Begründen}
\abhakgruppe{Einheit 3 · Verflechtung}
\abhak{13}{Verflechtung}
\abhak{14}{Verflechtung – Fehler finden, Begründen, Anwendung}
\abhakgruppe{Einheit 4 · Übergangsmodell}
\abhak{15}{Übergangsmodell}
\abhak{16}{Übergangsmodell – Fehler finden, Begründen, Anwendung}
\abhakgruppe{Einheit 5 · Stationär und langfristig}
\abhak{17}{Stationär und langfristig}
\abhak{18}{Übergangsprozess: Eignung eines Populationsmodells zur langfristigen Beschreibung beurteilen}
\abhak{19}{Stationär und langfristig – Fehler finden, Begründen, Anwendung}
\end{abhakseite}

\end{document}
